% part: intuitionistic-logic
% chapter: soundness-completeness
% section: canonical-model

\documentclass[../../../include/open-logic-section]{subfiles}

\begin{document}

\olfileid{int}{sc}{mod}

\olsection{典范模型}

我们模型中的世界是自然数的有穷序列~$\sigma$，即 $\sigma \in \Nat^*$。注意，$\Nat^*$ 由以下规则归纳定义：
\begin{enumerate}
\item $\emptyseq \in \Nat^*$。
\item 若 $\sigma \in \Nat^*$ 且 $n \in \Nat$，则 $\sigma.n \in \Nat^*$（其中 $\sigma.n$ 指 $\sigma \concat \tuple{n}$，而 $\sigma \concat \sigma'$ 是 $\sigma$ 与 $\sigma'$ 的拼接）。
\item 除此之外，没有其他元素属于 $\Nat^*$。
\end{enumerate}

因此，我们可以用 $\Nat^*$ 给出归纳定义。

令 $\tuple{!B_1, !C_1}$, $\tuple{!B_2, !C_2}$, \dots, 为所有公式对的一个枚举。给定一个公式集~$\Delta$，按如下方式归纳地定义 $\Delta(\sigma)$：
\begin{enumerate}
\item $\Delta(\emptyseq) = \Delta$
\item $\Delta(\sigma.n) = {}$
  \[
  \begin{cases}
    (\Delta(\sigma) \cup \{!B_n\})^* &
    \text{若 $\Delta(\sigma) \cup \{!B_n\} \Proves/ !C_n$} \\
    \Delta(\sigma) & \text{否则}
  \end{cases}
  \]
\end{enumerate}

这里的 $(\Delta(\sigma) \cup \{!B_n\})^*$ 是指将 \olref[lin]{lem:lindenbaum} 应用于集合 $\Delta(\sigma) \cup \{!B_n\}$ 与公式~$!C_n$ 所得到的公式质集。注意，根据此定义，如果 $\Delta(\sigma) \cup \{!B_n\} \Proves/ !C_n$，那么 $\Delta(\sigma.n)
\Proves !B_n$ 且 $\Delta(\sigma.n) \Proves/ !C_n$。此外注意，对任意~$n$，$\Delta(\sigma) \subseteq \Delta(\sigma.n)$。如果 $\Delta$ 是质集，那么对所有~$\sigma$，$\Delta(\sigma)$ 也是质集。

\begin{defn}\ollabel{defn:canonical-model}
  假设 $\Delta$ 是质集。则 $\Delta$ 的\emph{典范模型} $\mModel{M(\Delta)}$ 定义如下：
  \begin{enumerate}
  \item $W = \Nat^*$，即自然数的有穷序列的集合。
  \item $R$ 是偏序关系，其中 $R\sigma\sigma'$ 当且仅当 $\sigma$ 是~$\sigma'$ 的初始段（即存在序列~$\sigma''$ 使得 $\sigma' = \sigma \concat \sigma''$）。
  \item $V(p) = \Setabs{\sigma}{p \in \Delta(\sigma)}$。
  \end{enumerate}
\end{defn}

容易验证 $R$ 是一个偏序关系。同时，关于~$V$ 的单调性条件也成立。由 $\Delta(\sigma)
\subseteq \Delta(\sigma.n)$，通过对~$\sigma$ 进行归纳，可得每当 $R\sigma\sigma'$ 时，有 $\Delta(\sigma) \subseteq \Delta(\sigma')$。

\end{document}